\documentclass[11pt,twocolumn]{article}

% Packages
\usepackage[utf8]{inputenc}
\usepackage[T1]{fontenc}
\usepackage{lmodern}
\usepackage[margin=1in]{geometry}
\usepackage{graphicx}
\usepackage{amsmath}
\usepackage{amssymb}
\usepackage{booktabs}
\usepackage{hyperref}
\usepackage{xcolor}
\usepackage{float}
\usepackage{caption}
\usepackage{subcaption}
\usepackage{algorithm}
\usepackage{algpseudocode}
\usepackage{listings}
\usepackage{natbib}
\usepackage{fancyhdr}
\usepackage{titlesec}

% Hyperref setup
\hypersetup{
    colorlinks=true,
    linkcolor=blue!60!black,
    citecolor=blue!60!black,
    urlcolor=blue!60!black
}

% Header/Footer
\pagestyle{fancy}
\fancyhf{}
\fancyhead[L]{\small Sollertia: Fine Motor Control in XR}
\fancyhead[R]{\small \thepage}
\renewcommand{\headrulewidth}{0.4pt}

% Title formatting
\titleformat{\section}{\large\bfseries}{\thesection}{1em}{}
\titleformat{\subsection}{\normalsize\bfseries}{\thesubsection}{1em}{}

% Code listing style
\lstset{
    basicstyle=\ttfamily\small,
    breaklines=true,
    frame=single,
    numbers=left,
    numberstyle=\tiny\color{gray},
    keywordstyle=\color{blue},
    commentstyle=\color{green!50!black},
    stringstyle=\color{red!60!black}
}

\title{
    \vspace{-1cm}
    \textbf{Sollertia: Measuring Fine Motor Control in Extended Reality Through Task-Based Interaction and Wearable Force Sensing}
}

\author{
    Anaya Yorke\textsuperscript{1} \quad
    Garret S. Stand\textsuperscript{1} \\[0.5em]
    \textsuperscript{1}Independent Research \\[0.3em]
    \small\texttt{github.com/anaya33/Sollertia}
}

\date{\today}

\begin{document}

\maketitle

\begin{abstract}
We present Sollertia, a system for measuring fine motor control in extended reality (XR) using task-based interaction and wearable force input. The system replicates the same button-pressing task in both physical and XR conditions, enabling direct comparison of motor performance across environments. By combining behavioral metrics (reaction time, spatial accuracy, movement trajectory) with continuous force data from finger-mounted sensors, Sollertia captures how people move, react, and apply pressure during interaction. This paper describes the system architecture, measurement methodology, and current study design. The goal is to evaluate whether XR interaction can capture meaningful fine motor behavior comparable to physical tasks.
\end{abstract}

\section{Introduction}

Extended reality (XR) technologies are increasingly used for tasks requiring precise motor control, yet a fundamental question remains: \textit{Do motor behaviors observed in XR accurately reflect real-world motor capabilities?}

Traditional motor assessments provide limited quantitative insight, often yielding only summary scores that obscure the temporal dynamics of movement \citep{schoen2025pegs}. XR environments offer the potential for precise, objective measurement of motor behavior \citep{chen2024metrics}, but validation against physical task performance is needed.

This paper introduces Sollertia, a system designed to address this question through:

\begin{enumerate}
    \item \textbf{Standardized task replication} across physical and XR environments
    \item \textbf{Multi-modal measurement} combining behavioral and force metrics
    \item \textbf{High temporal resolution} data capture for movement dynamics analysis
\end{enumerate}

\section{Background}

\subsection{Motor Control Assessment}

Traditional clinical assessments of fine motor control rely on timed tasks that yield single summary scores (e.g., time to complete). While clinically validated, these measures obscure the temporal dynamics of movement—including reaction time, movement trajectory, and force modulation—that may provide earlier or more sensitive indicators of motor function \citep{wei2019tapping}.

\subsection{Fitts' Law and Pointing Tasks}

Fitts' Law \citep{fitts1954information} provides a foundational model for understanding speed-accuracy tradeoffs in aimed movements:

\begin{equation}
    MT = a + b \cdot ID
\end{equation}

where $MT$ is movement time, $ID = \log_2(2D/W)$ is the index of difficulty, $D$ is movement distance, and $W$ is target width. This relationship has been validated extensively in 2D interfaces and, more recently, in 3D virtual environments \citep{lu2019modeling}.

\subsection{XR for Motor Assessment}

XR environments offer several advantages for motor assessment:

\begin{itemize}
    \item Precise control over task parameters
    \item Automatic, objective data collection
    \item Repeatable, standardized conditions
    \item Rich spatial and temporal data capture
\end{itemize}

However, validation that XR-based assessments correlate with physical task performance remains limited \citep{schneider2021accuracy}.

\section{System Design}

Sollertia comprises three integrated components: an XR application, wearable hardware, and a clinical dashboard.

\subsection{XR Application}

The Unity-based application presents a discrete pointing task on Meta Quest 3 with hand tracking. Visual targets appear on a virtual surface, and participants press them using natural finger movements.

\textbf{Technical specifications:}
\begin{itemize}
    \item Unity 6 with Universal Render Pipeline
    \item OpenXR + XR Interaction Toolkit
    \item Hand tracking (no controllers)
    \item WebSocket data streaming
    \item Millisecond-precision event logging
\end{itemize}

\subsection{Wearable Hardware}

Force-sensing resistors (FSRs) mounted on the index and middle fingers capture pressure data during interaction.

\begin{table}[H]
\centering
\caption{Hardware Specifications}
\begin{tabular}{ll}
\toprule
\textbf{Component} & \textbf{Details} \\
\midrule
Microcontroller & ELEGOO UNO R3 \\
Sensors & 2$\times$ FSR402 \\
Sampling Rate & 10 Hz \\
Communication & USB Serial (9600 baud) \\
Resolution & 10-bit ADC (0--1023) \\
\bottomrule
\end{tabular}
\label{tab:hardware}
\end{table}

The index and middle fingers were selected as primary measurement points due to their functional importance in precision grip and established clinical significance in motor assessment.

\subsection{Dashboard}

A Rust-based application aggregates data streams, providing:

\begin{itemize}
    \item Real-time visualization of force and events
    \item Session configuration and control
    \item Data export for offline analysis
\end{itemize}

\section{Measurement Framework}

\subsection{Behavioral Metrics}

The following metrics are computed from XR application data:

\begin{table}[H]
\centering
\caption{Behavioral Metrics}
\begin{tabular}{lp{4cm}}
\toprule
\textbf{Metric} & \textbf{Description} \\
\midrule
Reaction Time & Target onset to movement initiation \\
Movement Time & Movement start to target contact \\
Spatial Error & Distance from contact to target center \\
Trajectory & 3D path during reach \\
Variability & SD across trials \\
\bottomrule
\end{tabular}
\label{tab:behavioral}
\end{table}

\subsection{Force Metrics}

FSR data enables characterization of force application:

\begin{table}[H]
\centering
\caption{Force Metrics}
\begin{tabular}{lp{4cm}}
\toprule
\textbf{Metric} & \textbf{Description} \\
\midrule
Peak Force & Maximum pressure during press \\
Force Onset & Time to reach force threshold \\
Force Duration & Time above threshold \\
Force CV & Coefficient of variation \\
\bottomrule
\end{tabular}
\label{tab:force}
\end{table}

\subsection{Derived Metrics}

Composite metrics provide summary characterizations:

\begin{equation}
    \text{Throughput} = \frac{ID}{MT}
\end{equation}

\begin{equation}
    \text{Accuracy Index} = 1 - \frac{\text{error}}{\text{target radius}}
\end{equation}

\begin{equation}
    \text{Force Stability} = \frac{1}{CV_{\text{force}}}
\end{equation}

\section{Task Protocol}

\subsection{Discrete Pointing Task}

Each trial follows a structured sequence:

\begin{enumerate}
    \item \textbf{Fixation} (500 ms): Brief pause
    \item \textbf{Target Onset}: Visual target illuminates
    \item \textbf{Reach}: Participant moves toward target
    \item \textbf{Contact}: Fingertip contacts surface
    \item \textbf{Press}: Force applied until threshold
    \item \textbf{Feedback}: Confirmation of success
\end{enumerate}

\subsection{Task Parameters}

\begin{table}[H]
\centering
\caption{Configurable Task Parameters}
\begin{tabular}{lcc}
\toprule
\textbf{Parameter} & \textbf{Default} & \textbf{Range} \\
\midrule
Session Duration & 45 s & 30--120 s \\
Target Count & 9 & 4--16 \\
Target Diameter & 40 mm & 20--60 mm \\
Force Threshold & 1.5 N & 0.5--3.0 N \\
\bottomrule
\end{tabular}
\label{tab:params}
\end{table}

\section{Data Collection}

Each trial generates a synchronized record containing:

\begin{itemize}
    \item Trial and target identifiers
    \item Timestamps: stimulus, movement onset, contact, press, release
    \item Spatial data: contact position, target position, trajectory
    \item Force signal: continuous FSR readings
\end{itemize}

Data synchronization is achieved through timestamp alignment between the XR application and hardware streams.

\section{Analysis Pipeline}

\subsection{Preprocessing}

\begin{enumerate}
    \item Timestamp alignment across data streams
    \item Outlier detection and removal
    \item Signal filtering (low-pass for force data)
\end{enumerate}

\subsection{Feature Extraction}

Per-trial computation of behavioral and force metrics as defined in Section 4.

\subsection{Aggregation}

Session-level statistics:
\begin{itemize}
    \item Central tendency (mean, median)
    \item Dispersion (SD, IQR)
    \item Trends across trials (learning effects)
\end{itemize}

\subsection{Condition Comparison}

Statistical comparison of physical vs. XR performance using paired analyses.

\section{Validation Approach}

\subsection{Research Questions}

\begin{enumerate}
    \item Can XR-based motor tasks capture meaningful fine motor behavior?
    \item How does motor performance in XR compare to equivalent physical tasks?
    \item What role does force modulation play in characterizing motor control?
\end{enumerate}

\subsection{Planned Studies}

\textbf{Study 1: Test-Retest Reliability}

Assess consistency of metrics across repeated sessions within participants.

\textbf{Study 2: Physical-XR Comparison}

Compare performance on matched tasks in physical and XR conditions.

\textbf{Study 3: Cross-Environment Correlation}

Correlate metrics between physical and XR conditions to assess measurement validity.

\section{Discussion}

\subsection{Contributions}

Sollertia provides a framework for rigorous comparison of motor performance across physical and XR environments. Key contributions include:

\begin{itemize}
    \item Integration of behavioral and force metrics
    \item Standardized task protocol for cross-environment comparison
    \item Open-source implementation for research reproducibility
\end{itemize}

\subsection{Limitations}

Current limitations include:
\begin{itemize}
    \item Limited sensor coverage (two fingers)
    \item Moderate force sampling rate (10 Hz)
    \item Preliminary validation status
\end{itemize}

\subsection{Future Directions}

\begin{itemize}
    \item Additional sensor modalities (IMU, EMG)
    \item Machine learning for motor pattern classification
    \item Clinical validation studies
    \item Adaptive difficulty algorithms
\end{itemize}

\section{Conclusion}

Sollertia represents a step toward rigorous, quantitative measurement of fine motor control in XR environments. By combining behavioral and force metrics in a standardized task framework, the system enables direct comparison of motor performance across physical and virtual conditions.

The open-source implementation is available at \url{https://github.com/anaya33/Sollertia}.

\section*{Acknowledgments}

This project originated at UGAHacks XI and continues as an independent research effort. We thank the hackathon organizers and mentors for their support.

\bibliographystyle{plainnat}
\bibliography{references}

\end{document}
